\documentclass[aspectratio=169,11pt]{beamer}

\usetheme{Madrid}
\usecolortheme{default}
\usefonttheme{professionalfonts}
\setbeamertemplate{navigation symbols}{}
\setbeamertemplate{footline}[frame number]

\usepackage{amsmath, amssymb, booktabs}

\title{Space And Labor Markets}
\subtitle{Urban Access, Commuting, Amenities, And Local Opportunity}
\author{Special Topic 4: Urban and Labor -- Week 1}
\date{}

\begin{document}

\begin{frame}
  \titlepage
\end{frame}

\begin{frame}{Central question}
\begin{itemize}
    \item How does space enter labor-market choice, wage determination, and worker welfare?
    \item Workers choose over job-residence bundles, not jobs from a placeless menu.
    \item Wages, rents, commuting costs, amenities, risks, and outside options jointly define opportunity.
    \item Week 1 discipline: separate workplace, residence, access, wages, rents, and welfare.
\end{itemize}
\end{frame}

\begin{frame}{Why this is labor economics}
\begin{itemize}
    \item Space changes which jobs workers can search for, reach, accept, and keep.
    \item Local labor markets are shaped by commuting and search frictions, not only by administrative borders.
    \item A high nominal wage can be offset by rent, travel time, congestion, risk, or weak outside options.
    \item Urban wage premia raise labor questions: productivity, learning, sorting, or compensation?
\end{itemize}
\end{frame}

\begin{frame}{Worker job-residence choice}
\[
U_{ijr} = w_j - R_r - \tau(d_{jr}) + A_r + \varepsilon_{ijr}
\]
\small
\begin{tabular}{p{0.25\linewidth} p{0.60\linewidth}}
\toprule
Term & Interpretation \\
\midrule
$w_j$ & earnings at workplace $j$ \\
$R_r$ & housing cost at residence $r$ \\
$\tau(d_{jr})$ & generalized commuting cost between $j$ and $r$ \\
$A_r$ & amenity or disamenity value at residence $r$ \\
$\varepsilon_{ijr}$ & worker-specific fit with the bundle \\
\bottomrule
\end{tabular}
\end{frame}

\begin{frame}{The city as a labor-market bundle}
\small
\begin{tabular}{p{0.25\linewidth} p{0.60\linewidth}}
\toprule
Object & Labor meaning \\
\midrule
Workplace wage & pay attached to a job location \\
Residential rent & cost of accessing a place while working \\
Commuting cost & time, money, reliability, effort, and risk \\
Amenities & local nonwage welfare components \\
Job access & jobs reachable within a travel-cost budget \\
Outside option & best connected alternative nearby or elsewhere \\
\bottomrule
\end{tabular}
\vspace{0.35cm}
\begin{itemize}
    \item The city is not a label. It is a set of constraints and opportunities.
\end{itemize}
\end{frame}

\begin{frame}{Local labor markets and spatial search}
\[
\mathcal{J}_i(c) = \{j : \tau(d_{ij}) \leq c\}
\]
\begin{itemize}
    \item Accessible jobs decline as distance, time, cost, congestion, unreliability, or risk rises.
    \item Access is the feasible set; realized commute is the chosen match.
    \item Flow-based markets use commuting, travel time, and search links.
    \item Administrative markets are useful containers, but not automatically the economic market.
\end{itemize}
\end{frame}

\begin{frame}{Residence and workplace are different objects}
\small
\begin{tabular}{p{0.28\linewidth} p{0.56\linewidth}}
\toprule
Measure & Main question \\
\midrule
Residence-based wages & What do residents earn and face? \\
Workplace-based wages & What jobs and pay are located here? \\
Commuting flows & How are residences connected to workplaces? \\
Local rents & What cost is capitalized into residential access? \\
Travel-time markets & Which jobs are actually reachable? \\
\bottomrule
\end{tabular}
\vspace{0.35cm}
\begin{itemize}
    \item Conflating residence and workplace can misstate welfare, policy incidence, and local adjustment.
\end{itemize}
\end{frame}

\begin{frame}{Commuting costs, rents, and amenities}
\begin{itemize}
    \item A job offer is valuable only if it can be paired with a feasible residence.
    \item A residence is valuable partly because of the jobs it gives access to.
    \item Commuting links labor demand to households without requiring migration.
    \item Amenities and disamenities enter welfare even when payroll data omit them.
    \item Housing and transportation constraints can make the same wage worth different amounts to different workers.
\end{itemize}
\end{frame}

\begin{frame}{Spatial equilibrium and welfare}
\[
w_\ell - R_\ell + A_\ell = \bar{U}
\]
\begin{itemize}
    \item Wages, rents, amenities, and mobility jointly allocate workers across places.
    \item High wages need not imply high utility when rents or commuting costs are high.
    \item Low wages need not imply low utility when housing is cheap or amenities are valuable.
    \item When mobility is constrained, equilibrium gaps become labor-market objects.
\end{itemize}
\end{frame}

\begin{frame}{Urban wage premium}
\[
\Delta w^{urban}
= \Delta w^{productivity}
+ \Delta w^{learning}
+ \Delta w^{sorting}
+ \Delta w^{compensating}
\]
\small
\begin{tabular}{p{0.26\linewidth} p{0.58\linewidth}}
\toprule
Channel & Labor interpretation \\
\midrule
Productivity & density raises output or match quality \\
Learning & workers accumulate skills faster in cities \\
Sorting & higher-productivity workers select into dense places \\
Compensation & wages offset rents, congestion, risk, or disamenities \\
\bottomrule
\end{tabular}
\end{frame}

\begin{frame}{Job access and the geography of work}
\begin{itemize}
    \item Jobs can move across space even when workers do not.
    \item Job suburbanization changes who can reach employment at acceptable commuting cost.
    \item Access shocks can affect employment without a simple wage-price story.
    \item The empirical object is not only distance; it is generalized travel cost, information, schedules, safety, and feasible job matches.
\end{itemize}
\end{frame}

\begin{frame}{Methods discipline}
\begin{itemize}
    \item Separate residence-based and workplace-based measurement.
    \item Separate nominal wages from real opportunity.
    \item Separate access from realized commute.
    \item Separate spatial selection from causal local effects.
    \item Separate geographic boundaries from flow-based labor markets.
\end{itemize}
\end{frame}

\begin{frame}{Week 1 lab}
\begin{itemize}
    \item Reproduce: compute job access in a deterministic synthetic city.
    \item Diagnose: classify access, commute, wage, rent, amenity, and equilibrium objects.
    \item Transfer: study a small commuting shock and ask who adjusts through commuting versus migration.
    \item Goal: turn a place-specific setting into a portable labor-market mechanism.
\end{itemize}
\end{frame}

\begin{frame}{Bridge to Week 2}
\begin{itemize}
    \item Week 1 defines space as jobs, residences, commuting, rents, amenities, risks, and outside options.
    \item Week 2 asks when housing and moving costs block workers from reaching better labor-market opportunities.
    \item Central next question: who captures a local wage gain when rents and mobility constraints adjust?
\end{itemize}
\end{frame}

\begin{frame}{Takeaways}
\begin{itemize}
    \item Local labor markets are access systems, not just places on a map.
    \item Workplace and residence must be separated before interpreting wages or welfare.
    \item Urban wage premia are labor-market objects with multiple possible channels.
    \item The rest of the course studies constraints that reshape the feasible set for work.
\end{itemize}
\end{frame}

\end{document}
